% =====================================================================
%  CHAPTER 2 — THEORETICAL BACKGROUND
%  Rules that hold here (Ch. 1 \S1.6):
%   - a primary source for every bound and every definition
%   - no verbatim quotation from a paper that is not in the corpus/
%   - the framework of validity does NOT go here. It belongs to the
%     methodology chapter, as \S1.4 promises.
%
%  ENGLISH IS THE MASTER. content/el/ch2.tex follows this file, block
%  for block.
% =====================================================================

%%blk background.head
\chapter{Theoretical background: models, bounds and criteria of
correctness}
\label{ch:background}

%%blk background.intro.01
Chapter~\ref{ch:intro} showed that every consensus protocol is a
position in a design space. This chapter defines the axes of that space.
Without them, the comparisons of the later chapters have no units.

%%blk background.intro.02
The order is deliberate. First comes the environment in which a protocol
runs, that is the synchrony and the faults. Then comes the problem that
it solves. Then come the bounds that no protocol breaks. Last come the
criteria that judge its correctness and its performance.

%%blk background.synchrony.head
\section{Two axes of synchrony, not one}
\label{sec:synchrony}

%%blk background.synchrony.01
The literature often speaks of \enquote{synchronous} and
\enquote{asynchronous} systems as if one axis existed. Dwork, Lynch and
Stockmeyer separate two \cite{dwork1988consensus}:

%%blk background.synchrony.axes
\begin{enumerate}[label=\textbf{(\alph*)},leftmargin=2.4em]
  \item \textbf{The communication.} Is there an upper bound $\Delta$ on
  the delivery time of a message?
  \item \textbf{The processors.} Is there an upper bound $\Phi$ on the
  rate at which the clock of one node runs faster than the clock of
  another?
\end{enumerate}

%%blk background.synchrony.02
A system is \emph{synchronous} when both bounds exist and are known in
advance. It is \emph{asynchronous} when \emph{neither of the two}
exists. The two axes are independent, so mixed models also exist between
them, where one dimension is synchronous and the other is not. The
separation is not pedantry. For exactly this reason the same authors
give a table of resilience with five columns instead of two, and
Table~\ref{tab:dls} reproduces three of them.

%%blk background.synchrony.sub.partial
\subsection{The two forms of partial synchrony}

%%blk background.synchrony.03
Between the two extremes lies \emph{partial synchrony}. The same paper
defines it in two ways, and the difference between them matters for this
thesis.

%%blk background.synchrony.04
In the \textbf{first form} the bound $\Delta$ exists, but we do not know
it. We ask for a protocol that works correctly whatever the value is,
that is a protocol with no $\Delta$ built into it.

%%blk background.synchrony.05
In the \textbf{second form} we know $\Delta$, but the message system
does not always respect it. The authors add one constraint. In every
execution there is a moment in time that the nodes do not know. From
that moment on, the message system respects $\Delta$. They call it the
\emph{global stabilization time} and they abbreviate it GST.

%%blk background.synchrony.06
The same two forms are defined for $\Phi$ as well. For that reason the
table of the paper has a separate column for partially synchronous
processors with synchronous communication. This thesis uses the second
form for both axes, and so do the protocols that it evaluates.

%%blk background.synchrony.fig.synchrony
\dgfigure{fig:synchrony}%
{The three synchrony models, with one vertical stroke per message.}%
{The three synchrony models, with one vertical stroke per message. In
the synchronous model every message arrives inside a known $\Delta$. In
the asynchronous model there is no bound. In the partially synchronous
model there is an unknown moment GST, after which $\Delta$ holds
\cite{dwork1988consensus}. The GST is a \textbf{model of analysis} and
not a mechanism. No code computes it. The term belongs to the paper of
1988 itself and is not a later addition.}%
{% fig-synchrony — the three timing models, one vertical stroke per message.
% Structure figure: dgflow. dgCFT marks a message that arrives inside a known
% bound, dgBFT one that does not — the standard's own key/bad pairing.
% The axis is 7.5 cm long so that the whole figure, row labels included, comes
% out at \dgwidth like every other figure; the message positions are the
% original ones scaled by 7.5/9.6.
\begin{dgflow}
  % --- synchronous -------------------------------------------------------
  \draw[dgaxis] (0,3.0) -- (7.5,3.0);
  \node[anchor=east, inner sep=4pt] at (-0.1,3.0) {\fl{sync.model.synchronous}};
  \node[dgnote, anchor=north west] at (0,2.84) {\fl{sync.note.synchronous}};
  \foreach \x in {0.781,1.875,2.969,4.063,5.156,6.250}{
    \draw[dgarrow,-,draw=dgCFT] (\x,2.88) -- (\x,3.12);}

  % --- asynchronous ------------------------------------------------------
  \draw[dgaxis] (0,1.6) -- (7.5,1.6);
  \node[anchor=east, inner sep=4pt] at (-0.1,1.6) {\fl{sync.model.asynchronous}};
  \node[dgnote, anchor=north west] at (0,1.44) {\fl{sync.note.asynchronous}};
  \foreach \x in {0.469,0.859,2.344,5.078,5.313,6.719}{
    \draw[dgarrow,-,draw=dgBFT] (\x,1.48) -- (\x,1.72);}

  % --- partially synchronous --------------------------------------------
  \draw[dgaxis] (0,0.2) -- (7.5,0.2);
  \node[anchor=east, inner sep=4pt] at (-0.1,0.2) {\fl{sync.model.partial}};
  \foreach \x in {0.469,1.172,2.031,2.891}{
    \draw[dgarrow,-,draw=dgBFT] (\x,0.08) -- (\x,0.32);}
  \foreach \x in {4.063,4.922,5.781,6.641}{
    \draw[dgarrow,-,draw=dgCFT] (\x,0.08) -- (\x,0.32);}
  % The dashed line stops above the note: in the original it ran through it.
  \draw[dgaxis,-,dashed] (3.594,-0.45) -- (3.594,0.85);
  \node[dgtag, fill=dgPaper, text=dgAux, anchor=south] at (3.594,0.85) {\textbf{GST}};
  \node[dgtag, fill=dgPaper, text=dgBFT, anchor=north] at (1.62,-0.05)
       {\fl{sync.liveness.none}};
  \node[dgtag, fill=dgPaper, text=dgCFT, anchor=north] at (5.60,-0.05)
       {$\le\Delta$, \fl{sync.liveness.guaranteed}};
  \node[dgnote, anchor=north west] at (-1.2,-0.75) {\fl{sync.note.safety}};
\end{dgflow}
}

%%blk background.synchrony.07
The second form is the one that the Byzantine protocols of this thesis
adopt \cite{yin2019hotstuff} and \cite{malkhi2023hotstuff2}\prep. The
practical consequence appears in Figure~\ref{fig:synchrony}. Safety
holds along the whole axis. Liveness is guaranteed only to the right of
the GST.

%%blk background.faults.head
\section{Fault models}
\label{sec:faults}

%%blk background.faults.01
The second axis asks what a broken node is allowed to do. This thesis
uses four grades, with rising malice.

%%blk background.faults.02
One convention of symbols, stated once for the whole chapter. The
letters $t$ and $f$ name the \emph{same} quantity, the maximum number of
faulty nodes. The chapter keeps the symbol of each primary source, so
that the values stay recognizable in it.

%%blk background.faults.list
\begin{description}[leftmargin=2.2em,style=nextline,font=\normalfont\bfseries]
  \item[Crash.] The node stops and never speaks again. The fault is an
  \emph{absence} of information. The more specific \emph{fail-stop}
  model adds that the other nodes detect the stop \emph{reliably}. The
  distinction returns in Section~\ref{sec:smr}, where the bound $t{+}1$
  of Schneider assumes fail-stop and not plain crash.
  \item[Omission.] The node loses messages, but everything that it sends
  is correct.
  \item[Byzantine with authentication.] The node behaves in an arbitrary
  way, but it cannot forge the signature of another node.
  \item[Byzantine without authentication.] No guarantee. The node says
  whatever it wants to whoever it wants, and it coordinates with the
  other faulty nodes \cite{lamport1982byzantine}.
\end{description}

%%blk background.faults.03
The behaviour that separates the Byzantine model from the others has a
name. It is \emph{equivocation}. The same node says different things to
different receivers, and the receivers cannot detect this on their own.
Figure~\ref{fig:equivocation} shows it.

%%blk background.faults.fig.equivocation
\dgfigure{fig:equivocation}%
{Crash against equivocation: the one behaviour that makes the Byzantine
model expensive.}%
{Crash against equivocation: the one behaviour that makes the Byzantine
model expensive. On the left, $K_2$ stops. $K_1$ and $K_3$ have less
information, but not contradictory information. On the right, $K_2$
sends a different value to each one and continues to work. Neither $K_1$
nor $K_3$ can detect this alone. Trusted hardware lowers the bound from
$3f{+}1$ to $2f{+}1$ exactly because it stops this behaviour. The
classification comes from the survey \cite{amiri2024bedrock} (App.~A,
O7). The primary source of the bound is the line of work on trusted
hardware to which that survey points: TrInc, which it names, and others
that it cites by reference number without naming them. This thesis has
\emph{not} checked them independently.}%
{% fig-equivocation — crash versus equivocation. Structure figure: dgflow.
% Left panel is the crash model, so every mark in it is dgCFT; right panel is
% the Byzantine model, so every mark in it is dgBFT. K1/K3 are ordinary nodes
% in both panels and stay dgbox grey. The two panels are 6 cm apart so that
% the pair comes out at \dgwidth.
\begin{dgflow}
  % ---- left: crash ------------------------------------------------------
  \node[text=dgCFT] at (1.5,3.35) {\textbf{\fl{equivocation.crash.title}}};
  \node[dgboxkey] (c2) at (0.0,1.2) {$K_2$};
  \node[dgbox]    (c1) at (3.0,2.4) {$K_1$};
  \node[dgbox]    (c3) at (3.0,0.0) {$K_3$};
  \draw[dgarrowkey] (c2) -- (c1)
        node[dgtag, fill=dgPaper, sloped, above, pos=0.5]
        {\fl{quote.open}$v=5$\fl{quote.close}};
  \draw[dgarrowkey] (c2) -- (c3)
        node[dgtag, fill=dgPaper, sloped, below, pos=0.5]
        {\fl{quote.open}$v=5$\fl{quote.close}};
  % K2 is gone: the cross sits just under the box it kills.
  \draw[dgarrowkey,-] (-0.30,0.65) -- (0.30,0.05);
  \draw[dgarrowkey,-] (-0.30,0.05) -- (0.30,0.65);
  \node[dgnote, anchor=north, align=center] at (1.5,-0.55)
       {\fl{equivocation.crash.note}};

  % ---- right: Byzantine -------------------------------------------------
  \node[text=dgBFT] at (7.5,3.35) {\textbf{\fl{equivocation.byzantine.title}}};
  \node[dgboxbad] (b2) at (6.0,1.2) {$K_2$};
  \node[dgbox]    (b1) at (9.0,2.4) {$K_1$};
  \node[dgbox]    (b3) at (9.0,0.0) {$K_3$};
  \draw[dgarrowbad] (b2) -- (b1)
        node[dgtag, fill=dgPaper, sloped, above, pos=0.5]
        {\fl{quote.open}$v=5$\fl{quote.close}};
  \draw[dgarrowbad] (b2) -- (b3)
        node[dgtag, fill=dgPaper, sloped, below, pos=0.5]
        {\fl{quote.open}$v=7$\fl{quote.close}};
  \node[dgnote, anchor=north, align=center] at (7.5,-0.55)
       {\fl{equivocation.byzantine.note}};
\end{dgflow}
}

%%blk background.faults.04
Table~\ref{tab:dls} combines the two axes. It is the most compressed
result of the chapter, and it explains why partial synchrony is not
free.

%%blk background.faults.tab.dls
\begin{table}[t]
\centering
\caption[The minimum number of nodes per model: the smallest $N$ for
which a consensus protocol that tolerates $t$ faults exists.]%
{The minimum number of nodes per model: the smallest $N$ for which a
consensus protocol that tolerates $t$ faults exists. The table
reproduces \textbf{three of the five columns} of Table~I of
\cite{dwork1988consensus}. The note justifies the choice. The symbol
$\infty$ means that no such $N$ exists, that is the FLP result.
\emph{Two readings that are worth the effort:} in the synchronous model
with authentication you tolerate \emph{any} number of faults, and
partial synchrony raises the cost of fail-stop from $t$ to $2t{+}1$,
that is you pay for not knowing when stability begins.}
\label{tab:dls}
\small
\begin{tabularx}{\linewidth}{@{}X c c c@{}}
\toprule
\textbf{Kind of fault} & \textbf{Synchronous} & \textbf{Asynchronous} &
\textbf{Partially synchronous} \\
\midrule
Fail-stop (in the sense of the source) & $t$ & $\infty$ & $2t+1$ \\
Omission                  & $t$      & $\infty$ & $2t+1$ \\
Byzantine with authentication & $t$  & $\infty$ & $3t+1$ \\
Byzantine                 & $3t+1$   & $\infty$ & $3t+1$ \\
\bottomrule
\end{tabularx}

\vspace{4pt}
\begin{minipage}{0.94\linewidth}\setstretch{1.08}\small
\textbf{Note on precision.} The column \enquote{Partially synchronous}
corresponds to the case where \emph{both} the communication \emph{and}
the processors are partially synchronous. The original gives five
columns, because it also separates the combinations where one dimension
is synchronous and the other is not. The version here keeps the three
that this thesis uses. The paper itself is outside the corpus. The
\emph{reference copy} is the open institutional copy of Lynch's group at
MIT (\texttt{groups.csail.mit.edu/tds/papers/Lynch/jacm88.pdf}). From
there we read the quantities of the table, the definitions of $\Delta$
and $\Phi$, the two forms of partial synchrony and the term GST. Nothing
is quoted verbatim.

\smallskip
\textbf{One word, two meanings.} The source names the first row
\emph{fail-stop}, and defines it as a processor that executes correctly
and can stop at any time, and once stopped cannot restart. There is
\textbf{no clause about detection}. That is not the sense in which
Section~\ref{sec:faults} of this chapter uses the word, following
\cite{schneider1990implementing}, where fail-stop adds that the other
nodes detect the stop \emph{reliably}. The difference decides the
asynchronous column: reliable detection is a strictly stronger
assumption than the one this table makes, and the $\infty$ in that column
belongs to the weaker assumption of the source. The row is therefore
labelled as the source labels it, and the reader is told which of the
two meanings is in force. This is the same trap as
Section~\ref{sec:smr}: one word, two counts.
\end{minipage}
\end{table}

%%blk background.smr.head
\section{Consensus, atomic broadcast and SMR are not the same}
\label{sec:smr}

%%blk background.smr.01
The three names are often used interchangeably. They are not synonyms,
and the confusion produces wrong comparisons in the later chapters.

%%blk background.smr.list
\begin{description}[leftmargin=2.2em,style=nextline,font=\normalfont\bfseries]
  \item[Consensus.] The nodes agree on \emph{one} value.
  \item[Atomic broadcast.] The nodes deliver the \emph{same sequence} of
  messages in the same order.
  \item[State machine replication (SMR).] The atomic broadcast is
  applied to a deterministic machine, so that all the replicas reach the
  same state \cite{schneider1990implementing}.
\end{description}

%%blk background.smr.02
Figure~\ref{fig:smr} summarizes the stack. SMR is the level where this
thesis lives. Schneider divides its requirement into two conditions.
Every non-faulty replica receives every request, and every non-faulty
replica processes the requests in the same relative order. The first
condition is called \emph{agreement}, the second one \emph{order}. Note
that the second one speaks of \emph{relative} order, and not of a global
timestamp.

%%blk background.smr.03
Above all of this stands an assumption that breaks easily in practice.
The sequence of requests must fix the outputs of the machine completely,
and nothing else must affect them \cite{schneider1990implementing}. If
the machine calls a generator of random numbers or a clock, SMR stops
holding.

%%blk background.smr.fig.smr
\dgfigure{fig:smr}%
{The stack: where each thing lives.}%
{The stack: where each thing lives. Consensus solves the order. The
state machine executes. All the protocols of this thesis live in the
middle box, and for that reason they compare. The determinism of the
machines is an assumption and not a conclusion
\cite{schneider1990implementing}.}%
{% fig-smr — where each layer lives. Structure figure: dgflow.
% The middle box is the subject of the whole thesis, so it is dgboxkey; the
% clients and the replicas are ordinary components, so they are dgbox. Both
% CFT and BFT protocols live in the middle box, which is why it cannot carry
% either series colour.
\begin{dgflow}
  \node[dgbox]    (cl) at (0.0,0.0) {\fl{smr.clients}};
  \node[dgboxkey] (co) at (4.0,0.0)
       {\bfseries\fl{smr.consensus}\\\mdseries\fl{smr.agreement.order}};
  \node[dgbox]    (sm1) at (8.4, 0.95) {\fl{smr.machine} 1};
  \node[dgbox]    (sm2) at (8.4, 0.00) {\fl{smr.machine} 2};
  \node[dgbox]    (sm3) at (8.4,-0.95) {\fl{smr.machine} 3};
  \draw[dgarrow] (cl) -- (co);
  \draw[dgarrow] (co) -- (sm1);
  \draw[dgarrow] (co) -- (sm2);
  \draw[dgarrow] (co) -- (sm3);
  \node[dgnote, anchor=north, align=center] at (8.4,-1.45)
       {\fl{smr.deterministic}\\\fl{smr.same.io}};
  \node[dgnote, anchor=north] at (4.0,-1.15)
       {\fl{smr.here.live} Paxos, Raft, PBFT, HotStuff};
\end{dgflow}
}

%%blk background.smr.sub.trap
\subsection{A counting trap: $2t{+}1$ and $3t{+}1$ do not count the same
thing}

%%blk background.smr.04
Schneider writes that Byzantine faults need at least $2t{+}1$ replicas,
while fail-stop faults need only $t{+}1$
\cite{schneider1990implementing}. The second value assumes that the
failure is \emph{detectable}. In that case it is enough to believe any
node that is alive. Chapter~\ref{ch:intro} writes $3f{+}1$ for the
Byzantine model. The two do not conflict, because they count different
things.

%%blk background.smr.trap.list
\begin{itemize}
  \item Schneider counts replicas whose \emph{outputs are voted on}.
  With $2t{+}1$ replicas the majority of the outputs stays correct even
  after $t$ faults. Agreement and order are already taken as solved.
  \item The bound $3f{+}1$ is the bound for \emph{solving} agreement,
  that is the bound for the consensus protocol itself
  \cite{pease1980reaching}.
\end{itemize}

%%blk background.smr.05
The two are easy to confuse, because both quantities are written as a
function of the number of faults. This thesis uses $3f{+}1$ everywhere,
because it measures consensus protocols and not the voting of outputs.

%%blk background.bounds.head
\section{What is impossible and what it costs}
\label{sec:bounds}

%%blk background.bounds.roadmap
Four results bound every protocol in this thesis. The first says that in
a fully asynchronous system no deterministic protocol always terminates.
The second gives the smallest number of nodes that tolerates a given
number of Byzantine faults. The third gives the cost in messages. The
fourth gives the smallest number of communication steps. They are taken
in that order, because each one narrows the space that the one before it
left open.

%%blk background.bounds.sub.flp
\subsection{FLP}

%%blk background.bounds.01
The fundamental negative result fits in one line. Fischer, Lynch and
Paterson prove that \enquote{No consensus protocol is totally correct in
spite of one fault} \cite{fischer1985impossibility}. The original is in
the corpus and the sentence is quoted verbatim.

%%blk background.bounds.02
The one line, however, carries two assumptions, and the paper itself
states both of them explicitly. The system is \emph{fully asynchronous}
and the protocol is \emph{deterministic}. The second assumption is the
way out. Randomized protocols do not fall under the theorem, and the
paper itself points to such work. All the protocols of this thesis are
deterministic, so the theorem binds them.

%%blk background.bounds.03
Three points of precision change the way the result is stated. We read
all three in the paper itself.

%%blk background.bounds.flp.list
\begin{enumerate}
  \item The paper does \textbf{not} define correctness through the named
  properties \enquote{agreement} and \enquote{validity}. The word
  \emph{validity} does not appear at all, and the word \emph{agreement}
  appears only as the name of the problem. It defines a protocol as
  \emph{partially correct} with two conditions. No accessible
  configuration has more than one decision value, and for every value
  there is an accessible configuration that decides it. The second
  condition is \emph{non-triviality}, which is weaker than validity.
  \item The paper does \textbf{not} examine Byzantine faults, and it
  states this explicitly. It assumes a reliable message system and one
  single crash. This makes the result \emph{stronger}. If consensus is
  impossible in the weaker fault model, it is all the more impossible in
  the Byzantine one.
  \item The requirement of termination is \textbf{extremely weak}, and
  the paper states this. It is enough that \emph{some} node decides
  eventually, and not all of them. The negative result therefore comes
  neither from malicious nodes nor from lost messages. It comes only
  from the fact that silence does not differ from delay.
\end{enumerate}

%%blk background.bounds.sub.third
\subsection{The bound of one third}

%%blk background.bounds.04
The bound $n \ge 3f+1$ for Byzantine agreement without authentication is
published first in 1980, in the form of the problem of \emph{interactive
consistency} \cite{pease1980reaching}. The name states the same
condition from the other side. The faulty nodes must be fewer than one
third of the nodes, so the correct ones are more than two thirds. Two
years later the same authors state it again as the allegory of the
Byzantine generals \cite{lamport1982byzantine}.

%%blk background.bounds.05
The same paper of 1980 already examines the \emph{authenticated} case as
well \cite{pease1980reaching}. There the faulty nodes can stay silent,
but they cannot alter what they relay. The bound of one third does not
hold in that case.

%%blk background.bounds.06
The allegory of 1982 adds the cleanest statement. The algorithm SM($m$)
solves the problem for any number of generals, provided that no one can
forge the signature of a loyal general \cite{lamport1982byzantine}. The
same paper adds one clause. The problem is empty with fewer than
$m{+}2$ generals. The result about signatures is therefore not a
contribution of 1982 but of 1980.

%%blk background.bounds.07
The bound $3f{+}1$ that all the protocols of this thesis use does
therefore not come from the absence of signatures. It comes from
\emph{partial synchrony}, as the row \enquote{Byzantine with
authentication} of Table~\ref{tab:dls} shows. There the requirement is
$n \ge 3t{+}1$ even when the messages are signed.

%%blk background.bounds.quote
\begin{quote}\small
If someone asks you \enquote{why do signatures not save you from
$3f{+}1$?}, the answer is that row of the table. Signatures are enough
in the \emph{synchronous} model. In the partially synchronous model they
are not enough.
\end{quote}

%%blk background.bounds.sub.messages
\subsection{The cost in messages}

%%blk background.bounds.08
Impossibility is not the only bound. There is also a lower bound on the
number of messages. Dolev and Reischuk prove that every
\textbf{deterministic} algorithm for Byzantine agreement that tolerates
$t$ faults has an execution in which the correct nodes send a number of
messages that is quadratic in $t$ \cite{dolev1985bounds}.

%%blk background.bounds.09
Three points need attention, and the chapter states all three here
instead of hiding them.

%%blk background.bounds.messages.list
\begin{itemize}
  \item The theorem is \textbf{not} limited to the unauthenticated
  algorithms. Its hypothesis contains no clause about authentication, so
  it covers both cases. Many secondary sources state this wrongly.
  \item The restriction to \textbf{determinism} is not decoration.
  Randomized protocols can reach subquadratic communication, and they
  pay with a probability of error. This does not hold always. Against a
  \emph{strongly adaptive} adversary, which erases messages after the
  fact, the quadratic bound holds for them as well
  \cite{abraham2019ccba}. They therefore do not break the bound. They
  step outside its hypothesis, and only under a weaker model of the
  adversary. All the protocols of this thesis are deterministic.
  \item The journal paper itself is behind a paywall and outside the
  corpus. We checked the content in the earlier conference version, and
  the numbering of the theorems may differ.
\end{itemize}

%%blk background.bounds.10
The consequence concerns the chapters on the Byzantine protocols
directly. When HotStuff advertises a \emph{linear} view change, the
claim holds for the view change as well, provided that the leader who
takes over is correct and the network is stable \cite{yin2019hotstuff}.
The lower bound is not violated. It moves to the worst case. The same
paper reports $O(fn)$ \emph{authenticators} over $f$ consecutive faulty
leaders, and that is $O(n^{2})$ when $f$ grows with $n$.

%%blk background.bounds.11
\textbf{Attention to the unit.} HotStuff counts \emph{authenticators},
while the bound of Dolev and Reischuk counts \emph{messages}. The two
quantities are not the same, and the HotStuff paper itself states this.
One message that carries $n$ signatures counts as much as $n$ messages
with one signature each. The comparison here is qualitative and not an
exact accounting.

%%blk background.bounds.sub.twostep
\subsection{The two-step bound}

%%blk background.bounds.12
The last category of bounds concerns the number of phases. \emph{Two
steps} means that a command commits with two consecutive phases only.
These are the proposal of the leader and one round of votes.

%%blk background.bounds.13
\textbf{Attention to the counting convention.} The survey
\cite{amiri2024bedrock} counts PBFT with \emph{three} phases of ordering
and FaB with \emph{two}. Table~\ref{tab:delays} follows the convention
of \cite{malkhi2023hotstuff2}\prep, where PBFT appears with two. The two
countings do not compare with each other.

%%blk background.bounds.14
Saving one phase cost $5f{+}1$ nodes at first \cite{martin2006fab}\unver.
With authentication $5f{-}1$ nodes are enough, and that is the lower
bound \cite{kuznetsov2021optimal,abraham2021goodcase}\unver.

%%blk background.bounds.15
The survey \cite{amiri2024bedrock} carries both quantities, and this
thesis took them from there. The text names the primary sources, but
this thesis has not checked them independently.
Table~\ref{tab:bounds} states this.

%%blk background.bounds.tab.bounds
\begin{table}[t]
\centering
\caption[Summary of the bounds that this thesis uses, with their primary
source.]%
{Summary of the bounds that this thesis uses, with their primary source.
The column \enquote{Documentation} is part of the evidence and has three
values. \emph{Corpus}: the paper exists and we read it.
\emph{Index}: the bibliographic data were verified in a DOI or in an
index, and the content in an open institutional copy, in an authentic
copy of the author, or in an earlier version. \emph{Unverified}: the
data come from the reference list of \cite{amiri2024bedrock}. Where the
column does not say \emph{corpus}, the text quotes nothing verbatim.}
\label{tab:bounds}
\small
\begin{tabular}{@{}l l c@{}}
\toprule
\textbf{Bound} & \textbf{Primary source} & \textbf{Documentation} \\
\midrule
Impossibility in asynchrony, 1 crash & \cite{fischer1985impossibility} & corpus \\
$n \ge 3f+1$, Byzantine & \cite{pease1980reaching} & index \\
Signatures: the bound collapses (synchronous) &
  \cite{pease1980reaching,lamport1982byzantine} & index \\
$n \ge 2f+1$ / $3f+1$ in partial synchrony & \cite{dwork1988consensus} & index \\
$\Omega(t^2)$ messages, deterministic & \cite{dolev1985bounds} & index \\
$5f{+}1$ for one phase fewer & \cite{martin2006fab}\unver\ & unverified \\
$5f{-}1$ two-step lower bound &
  \cite{kuznetsov2021optimal,abraham2021goodcase}\unver\ & unverified \\
\bottomrule
\end{tabular}
\end{table}

%%blk background.bounds.16
Table~\ref{tab:bounds} collects the bounds of the chapter together with
their documentation.

%%blk background.quorums.head
\section{Quorums: the one property that does the work}
\label{sec:quorums}

%%blk background.quorums.01
A \emph{quorum} is the smallest set of nodes that makes a decision
final. The only property that it needs is intersection.

%%blk background.quorums.sub.derivation
\subsection{The derivation}

%%blk background.quorums.02
Figure~\ref{fig:quorum} of Chapter~\ref{ch:intro} gives the sizes. Here
is the argument that produces them, in three steps.

%%blk background.quorums.list
\begin{enumerate}
  \item You cannot wait for more than $n-f$ answers. The $f$ faulty
  nodes may stay silent, and silence does not differ from delay. The
  quorum therefore has size $q = n-f$.
  \item Two quorums intersect in at least $2q - n = n - 2f$ nodes.
  \item In the Byzantine model the intersection must contain at least
  one \emph{honest} node. Up to $f$ Byzantine nodes hide inside the
  intersection, so $n - 2f > f$, that is $n \ge 3f+1$.
\end{enumerate}

%%blk background.quorums.03
In the crash model step~3 does not exist. It is enough that
$n - 2f > 0$, that is $n \ge 2f+1$. That one step is the whole
difference of cost between the two models.

%%blk background.quorums.quote
\begin{quote}\small
\textbf{What this is and what it is not.} The derivation above is an
argument about \emph{quorum design}. It shows that quorums of size $n-f$
need $n \ge 3f+1$. It is not the general proof of impossibility, which
belongs to \cite{pease1980reaching}.
\end{quote}

%%blk background.quorums.sub.intersection
\subsection{Intersection is required less than we thought}

%%blk background.quorums.04
From Paxos onward, the majority quorum was taken to be necessary.
Howard, Malkhi and Spiegelman show that it is not
\cite{howard2016flexiblepaxos}. In Paxos, intersection is required only
\emph{between} the two phases. The quorums of Phase~1 do not need to
intersect with each other, and neither do the quorums of Phase~2.

%%blk background.quorums.05
The practical consequence is that you can make the quorum of Phase~2
smaller, that is the quorum of the steady state. You pay with a larger
quorum in the election of the leader. The chapter on Paxos returns to
this, because it changes the reading of the arithmetic of
Figure~\ref{fig:quorum}.

%%blk background.correctness.head
\section{Criteria of correctness}
\label{sec:correctness}

%%blk background.correctness.roadmap
Two questions decide whether a replicated service is correct, and they
are separate questions. The first is what the protocol guarantees about
the order it builds; that is safety and liveness. The second is what a
client is entitled to observe from outside; that is linearizability.
This section takes them in that order, because the second is stated in
terms of the first.

%%blk background.correctness.sub.safety
\subsection{Safety and liveness}

%%blk background.correctness.01
Two categories of properties judge a protocol. \emph{Safety} says that
something bad never happens. \emph{Liveness} says that something good
happens eventually.

%%blk background.correctness.02
The asymmetry that governs all the protocols of this thesis is simple.
Safety holds at all times, even in a complete partition of the network.
Liveness is sacrificed. The choice is one of design and not an accident.
A system that stops can be repaired, while a system that wrote two
different values cannot.

%%blk background.correctness.sub.lin
\subsection{Linearizability}

%%blk background.correctness.03
Herlihy and Wing define the criterion that the client expects
\cite{herlihy1990linearizability}. Every operation appears to happen at
one instant between its invocation and its response. Operations that do
not overlap in time keep their order.

%%blk background.correctness.04
The second half does all the work. It is a constraint of \emph{real
time}, and for that reason linearizability costs coordination.

%%blk background.correctness.05
Two properties separate it from the serializability of the databases,
and both of them matter for this thesis.

%%blk background.correctness.list
\begin{description}[leftmargin=2.2em,style=nextline,font=\normalfont\bfseries]
  \item[Locality.] A system is linearizable if every object of it is
  linearizable. The property \emph{composes}. Neither serializability
  nor strict serializability is local
  \cite{herlihy1990linearizability}.
  \item[Nonblocking.] A pending invocation of a total operation is never
  forced to wait for another pending invocation. Serializability is
  inherently blocking \cite{herlihy1990linearizability}.
\end{description}

%%blk background.correctness.06
The same authors describe linearizability as a special case of strict
serializability, where each transaction is restricted to one operation
on one object.

%%blk background.performance.head
\section{What \enquote{performance} means for a protocol}
\label{sec:performance}

%%blk background.performance.01
This section is short and it carries weight. It defines, once, the
difference between two quantities that Chapter~\ref{ch:intro} showed to
be confused with each other.

%%blk background.performance.list
\begin{description}[leftmargin=2.2em,style=nextline,font=\normalfont\bfseries]
  \item[Phases.] How many consecutive rounds of proposal and vote a
  command needs before it commits. This is a property of the
  \emph{protocol} and it \emph{compares} between papers, provided that
  the counting convention is stated.
  \item[Message delays.] How many consecutive transmissions a command
  needs before it commits. In HotStuff one phase corresponds to two
  delays. Elsewhere the correspondence differs, and for that reason it
  is stated.
  \item[Real time.] How many milliseconds pass. This depends on the
  implementation, the hardware and the topology of the network. It does
  not compare between papers. This thesis reports its own measurements
  of real time in the evaluation chapter.
\end{description}

%%blk background.performance.02
Table~\ref{tab:delays} gives the first quantity, in \emph{phases}, and
its note states the conversion into delays. It is the only table of
performance in the chapter, because it is the only one that does not
depend on an environment.

%%blk background.performance.tab.delays
\begin{table}[t]
\centering
\caption[Phases on the happy path: how many phases one decision needs,
with a correct leader and after the GST.]%
{Phases on the happy path: how many phases one decision needs, with a
correct leader and after the GST. There are no milliseconds here, and
that is the point of the section. Sources:
\cite{malkhi2023hotstuff2}\prep\ (Table~1) for the four Byzantine
protocols, and \cite{ongaro2014raft} \S5.3 for the Multi-Paxos / Raft
row, where in the steady state \enquote{a new entry can be replicated
with a single round of RPCs to a majority of the cluster}.}
\label{tab:delays}
\small
\begin{tabular}{@{}l c c@{}}
\toprule
\textbf{Protocol} & \textbf{Phases} & \textbf{Model} \\
\midrule
Multi-Paxos / Raft & 1 & crash \\
PBFT               & 2 & Byzantine \\
Tendermint         & 2 & Byzantine, not responsive \\
HotStuff           & 3 & Byzantine \\
HotStuff-2\prep    & 2 & Byzantine \\
\bottomrule
\end{tabular}

\vspace{4pt}
\begin{minipage}{0.94\linewidth}\setstretch{1.08}\small
\textbf{Attention to the accounting.} The table counts \emph{phases},
with the same meaning and the same source that
Table~\ref{tab:viewchange} of Chapter~\ref{ch:intro} uses. In HotStuff
each phase corresponds to \emph{two} consecutive message delays, that is
the proposal of the leader and the votes. The source itself states this
\emph{for HotStuff}, so its three phases are six delays. The
correspondence is \textbf{not} universal. PBFT appears with $2$ because
it has two rounds of voting after the proposal, that is three delays in
total and not four. The table therefore compares \emph{phases as its
source counts them}, and not delays. The Multi-Paxos / Raft row does not
come from the same source and the caption marks it separately. The
number of phases does \textbf{not} translate into time without a
statement of the network. Three phases in a LAN cost less than two in a
wide-area network.
\end{minipage}
\end{table}

%%blk background.performance.03
The section closes with the rule that every later chapter applies. A
comparison of protocols compares phases. A comparison of implementations
compares times --- and only when the environment is stated. The full
framework of validity, the levels of comparability and the threats to
validity belong to the methodology chapter, as \S1.4 promises.
